
\textbf{Définition 2.1.}~: \uline{On appelle} pré-@@M0@@-anneau \uline{un anneau} @@M1@@, \uline{commutatif à élément unité}, \uline{muni d'un homomorphisme de groupes abéliens}~:
@@M2@@
\uline{tel que pour tout} @@M3@@, @@M4@@

On posera @@M5@@~; la donnée d'une structure de pré-@@M6@@-anneau sur un anneau @@M7@@ équivaut donc à la donnée d'une famille d'applications @@M8@@, @@M9@@, de @@M10@@ dans @@M11@@, telles que~:
@@M12@@

Soient @@M13@@ deux pré-@@M14@@-anneaux, et @@M15@@, un homomorphisme d'anneaux~; on dira que @@M16@@ est un homomorphisme de pré-@@M17@@-anneaux, ou encore un @@M18@@-\uline{homomorphisme}, si le diagramme~:
@@M19@@
est commutatif.

Ceci équivaut encore aux relations suivantes~:
@@M20@@

